\section{Risultati ottenuti}
L'obiettivo di implementare all'interno di FLY un supporto per i database MongoDB è stato raggiungo con successo. Esso ha portato alla nascita di un secondo obiettivo implementativo, ovvero offrire all'utente la possibilità di interagire con più istanze di database contemporaneamente. Per far fronte a questo secondo obiettivo, è stato implementato all'interno di FLY l'entità \verb|distributed-query| che permette appunto all'utente di effettuare una singola query su più database utilizzando una sintassi e un comportamento analogo a quello di una query eseguita su un singolo database. \\
Il lavoro svolto è stato diviso in 5 parti:
\begin{itemize}
    \item studio delle funzionalità già presenti all'interno di FLY che consentono all'utente di manipolare dati memorizzati su sorgenti esterne;
    \item studio delle tecnologie e metodologie necessarie per consentire all'utente di interagire con un database MongoDB;
    \item implementazione dell'entità \verb|nosql| partendo ed adattando la struttura utilizzata per istanziare un oggetto di tipo \verb|sql|;
    \item estensione del comportamento dell'entità \verb|query| su database MongoDB;
    \item implementazione dell'entità \verb|distributed-query| per permettere all'utente di effettuare una singola operazione su più cluster MongoDB.
\end{itemize}

\section{Sviluppi futuri}
I possibili sviluppi futuri che possono interessare il lavoro svolto durante questa tesi, e quindi l'ambito della gestione dei dati, sono:
\begin{itemize}
    \item l'estensione del comportamento dell'entità \verb|distributed-query| per i database MySQL;
    \item implementazione di un supporto ai database a grafo.
\end{itemize}

\subsection{Database a grafo}
I database a grafo sono una tipologia di database che utilizza nodi e archi per rappresentare e archiviare i dati. Questo modello di rappresentazione permette di ottenere una maggiore velocità nell'associazione di insiemi di dati rispetto ai database relazionali. Inoltre, scalano più facilmente a grandi quantità di dati e non richiedono le tipiche e onerose operazioni di unione. \\
Una possibile implementazione all'interno di FLY potrebbe essere quella di permettere all'utente di interagire con un database Neo4j. \cite{neo4j}. Neo4j è un Graph DBMS open source transazionale, prodotto dalla software house Neo Technology, sviluppato interamente in Java. Le sue proprietà principali sono:
\begin{itemize}
    \item transazioni ACID;
    \item può memorizzare miliardi di nodi e relazioni;
    \item alta velocità di interrogazione tramite attraversamenti;
    \item linguaggio di interrogazione dichiarativo e grafico;
    \item compatibile con i servizi di gestione ai database a grafo offerti da Azure e AWS.
\end{itemize}